\paragraph{Headline.} In the cascade, PPO beats every baseline, including
the rule with tuned thresholds. Against the tuned rule, the season-by-season
difference is $+1.79\pm0.67$: about 2.7 standard errors, with PPO better in
22 of the 30 test seasons (Table~\ref{tab:cascade_paired}). In the
single-reservoir problem the same comparison was a statistical tie. The
extra difficulty of coordinating two reservoirs around an uncontrolled
tributary is what lets the learned policy pull ahead. The size of the gain
over the tuned rule is modest (about 13\% of its penalty), but it is
consistent.

\paragraph{Where the gain comes from.} Table~\ref{tab:cascade_results}
splits the return into its parts. Compared with the tuned rule, PPO
spills 23\% less water (13{,}800 versus 17{,}900~m$^3$ per season), floods
the river on fewer days (8.5 versus 11.4), and has a slightly smaller
supply penalty with the same total shortfall ($\approx5\%$ of demand). The
three trajectory figures show the mechanisms:
\begin{itemize}
\item \textbf{Hedging in droughts (Figure~\ref{fig:cascade_traj_med}).}
In this median season a long dry spell empties both reservoirs under every
policy. The rules deliver the full demand until the water runs out and
then deliver almost nothing for about 20 days. PPO starts delivering
somewhat less than demand around day 15, \emph{before} the reservoirs are
empty, and stretches its storage further. Because the supply penalty is
\emph{squared}, many small shortfalls cost less than a few large ones:
spreading a 1{,}000~m$^3$/d deficit over twice as many days at half the
size halves its penalty. This is the classic \emph{hedging rule} of
reservoir operation theory. The agent was never told about it; it
discovered it from the shape of the reward.
\item \textbf{Anticipating storms (Figure~\ref{fig:cascade_traj_best}).}
In this storm-rich season PPO raises releases ahead of forecast storms, up
to and slightly above $\Qsafe$, and keeps more room in B, so the storms
around days 40--47 cause far less spill than under the rules.
\item \textbf{Where it loses (Figure~\ref{fig:cascade_traj_worst}).} In
the season where PPO does worst relative to the tuned rule, a cluster of
storms arrives at the very end of the season. The tuned rule happened to
be holding A very low (its 20\% threshold) and absorbed more of it. PPO,
which also weighs the end-of-season storage term, was fuller and spilled
more. A learned policy is better \emph{on average}, not in every season.
\end{itemize}

\paragraph{The rule baselines.} The threshold map
(Figure~\ref{fig:cascade_tuning}) shows the logic of the best rule: keep
the large upstream reservoir low ($\eta_A=0.2$) and the small downstream
one at half ($\eta_B=0.5$), so that A's slow, large storm volumes have room
while B buffers the fast local response. Even tuned, the rule cannot
\emph{change} these targets with the forecast or the season, and that
flexibility is what PPO adds. ``Supply only'', which never releases for
flood space, shows how much is at stake: it spills about 91{,}000~m$^3$ per
season and scores $-38.6$.

\paragraph{Lessons for applying RL.}
\begin{enumerate}\itemsep1pt
\item RL helps most where simple rules are structurally limited: several
interacting storages, uncontrolled inflows and forecasts. For one
reservoir, a well-tuned rule was as good.
\item The reward shapes the behavior. The hedging seen here follows from
the squared deficit penalty, and a linear penalty would not produce it.
\item Fair comparison requires tuned baselines, identical test seasons and
paired statistics. Against the untuned rule the gain would have looked
twice as large ($+3.7$).
\end{enumerate}
